% U1_regularity_v44.tex -- text v44 (fix27_texts_v44, 24 Nov 2026), built by fix27_texts_v44/tools/build_v44.py from fix24_texts_v43/U1_regularity_v43.tex (unchanged):
% the check fix25_check_v43/CHECK_v43.md (F-1..F-6, Note) and the LaTeX defects of fix26_latex_v43/LATEX_v43.md, where they touch this file.  Text v44 not yet compiled.
% U1_regularity_v44.tex -- text v43 (fix24_texts_v43, 23 Nov 2026), built by fix24_texts_v43/tools/build_v43.py from fix22_texts_v42/U1_regularity_v42.tex (unchanged):
% the findings of the delta check fix23_delta_check/DELTA_CHECK_v42.md (S-1..S-4, W-1..W-9), where they touch this file.  Not compiled.
% U1_regularity_v44.tex -- text v42 (fix22_texts_v42, 23 Nov 2026), built by fix22_texts_v42/tools/build_v42.py from fix21_texts_v41/U1_regularity_v41.tex (unchanged):
% Lean replaces DKS (writeup/lean_bridge/REVIEW_W7.md section 3) and the v41 leftovers, where they touch this file.  Not compiled.
% U1_regularity_v44.tex -- text v41 (fix21_texts_v41, 21 Nov 2026), built by fix21_texts_v41/tools/build_v41.py from fix19_texts_v40/U1_regularity_v40.tex (unchanged):
% wording items W-1..W-5 and the retag item of the delta check fix20_delta_check/DELTA_CHECK_v40.md where they touch this file.  Not compiled.
% U1_regularity_v44.tex -- text v40 (fix19_texts_v40, 21 Nov 2026), built by fix19_texts_v40/tools/build_v40.py from fix17_texts_v39/U1_regularity_v39.tex (unchanged):
% fixes of the final review fix18_final_review/FINAL_REVIEW_v39.md (F-1..F-7, R-1..R-3) where they touch this file.  Not compiled.
% U1_regularity_v44.tex -- text v39 (fix17_texts_v39, 20 Nov 2026), built by fix17_texts_v39/tools/build_v39.py from fix16_texts_v38/U1_regularity_v38.tex (unchanged):
% re-pinned step-4 records (manifest_v38 be909bbf...), v4 kits (WSL test runs only), step4_handoff_v44.tex; v38 retag fixes.  Not compiled.
% U1_regularity_v44.tex -- regularity lemma U-1 for step 3, text v38 (fix16_texts_v38; verifier fin13v4c, (A) chain v37, flux lemma v2).
% Built by tools/build_interface_v38.py from fix10_proofs/U1_regularity.tex (unchanged).  New in v38 (review fix15 D-1, D-5, C-5):
% X_S = 60; the open set O through GS-circ; B1 as Lambda <= L + 2C'; (S4.1)-(S4.3) by C-1 and GS-1; gap items 1 and 4 closed.
% Original header: regularity lemma U-1 (v4b verifier, (A) chain v36, flux lemma v2), with the modulation statement (P2').
% Fix round 10 (fix10_proofs), 20 Nov 2026.  Not compiled (no LaTeX in WSL).  Symbolic checks: U1_sym.py -> U1_sym.json (sympy, all_pass true).
% New file only; nothing closed is edited.
\documentclass[11pt]{article}
\usepackage[margin=1in]{geometry}
\usepackage{amsmath,amssymb,amsthm,booktabs,enumitem}
\newtheorem{lemma}{Lemma}\newtheorem{proposition}{Proposition}\newtheorem{remark}{Remark}\newtheorem{theorem}{Theorem}\newtheorem{definition}{Definition}
\newcommand{\fn}[1]{\texttt{\detokenize{#1}}}
\newcommand{\tH}{\textup{[hand]}}\newcommand{\tS}{\textup{[sym]}}\newcommand{\tR}{\textup{[read]}}\newcommand{\tL}{\textup{[Lean]}}\newcommand{\tI}{\textup{[inh]}}
\providecommand{\LeanSthree}{\textup{[Lean]}}% S3 items of review W7 (lean_bridge snap_s3); text v42 first marked them [hand; Lean formalization in progress]
\newcommand{\ip}[2]{\langle#1,#2\rangle}
\begin{document}
\title{Regularity lemma U-1 for step 3\\[2pt]\large the solution class, the modulation (P2$'$), and the energy identities at every time}
\date{24 November 2026 (text v44; text v43; text v42; text v41; text v40; text v39; first version fix round 10, reviewed in \fn{fix11_review} part C: sound modulo items). \emph{Text v43 compiled} (MiKTeX, \fn{fix26_latex_v43/LATEX_v43.md}); \emph{text v44 not yet compiled.}}
\maketitle

\noindent\textbf{Labels.} \tS: exact sympy identity in \fn{U1_sym.py} (this directory; \fn{U1_sym.json}, \fn{all_pass: true}); \tL: a compiled theorem of Micah Warren's Lean development \fn{CWFlow}, directory \fn{CWFlow/Curve/} or the library it imports (commit pending: files identified by the SHA256 sums of \fn{writeup/lean_bridge/snap_final/SHA256SUMS}; compiled per file against the repository's existing build outputs; axiom check \fn{[propext, Classical.choice, Quot.sound]}; \S\ref{U:sec:lean}; text v42: in text v41 \tL\ meant ``statement read from source, not compiled''); \LeanSthree: the S3 items of review W7 (S3a: periodic curve to \fn{Circle}; S3b: the shifted maximal flow and endpoint gluing), formalized after that review in \fn{Periodic2Circle.lean} and \fn{MaxShift.lean} (snapshot \fn{writeup/lean_bridge/snap_s3/SHA256SUMS}), printed as [Lean]; one macro, kept so that the provenance stays visible; the label [DKS] of text v41 is no longer used; \tR: verified by reading the cited file; \tI: inherited from a cited document, with the label it has there; \tH: hand argument written here, not machine-checked.

\noindent\textbf{Where U-1 is asked for.} \fn{fix5/ei_spec/EI_SPEC.md} \S3 (``Regularity \dots needs a short lemma, U-1'') and \S11; \fn{step3nl_v4c/V4C_REPORT.md} (open items); \fn{fix7_flux2/FLUX2_DESIGN.md} (table row ``energy identity, level 2'' and condition C8); \fn{fix8_flux2_impl/FLUX2_IMPL_REPORT.md} \S2 C8 and \S4.1: ``$\Lambda$ absolutely continuous; $g''''(\tau,\cdot)\in L^2(\mathrm{OR})$ a.e.; the level-2 identity with a $C^3$ cutoff; Agmon for $g''$ a.e.'' The continuation text \fn{fix1/b12_text/continuation_v33.tex} asks for (P1$'$)--(P3$'$), and \fn{P1_CONFIRMATION.md} \S4 item~2 asks for one sentence on (P2$'$): the modulation does not break down before $T_{\max}$ while the bootstrap holds.

\noindent\textbf{Summary.} Theorem~\ref{U:thm} gives all of these, in a stronger form: on any interval on which the modulation and the chart and box hypotheses hold, the Lyapunov functional and every energy used by the verifier, by Lemma~S(A) and by the flux lemma are $C^1$ in $\tau$, and their identities hold at \emph{every} $\tau$ (not only a.e.). Proposition~\ref{U:prop:mod} is (P2$'$). Inputs that are not proved here are listed in \S\ref{U:sec:inputs} and \S\ref{U:sec:gaps}.

%======================================================================
\section{Inputs, stated precisely}\label{U:sec:inputs}
\subsection{From step 4 (hand-off)}
\begin{enumerate}[label=(S4.\arabic*),leftmargin=3.2em]
\item \emph{(H6)} $\gamma(t_1):\mathbb S^1\to\mathbb C$ is a smooth immersion: Theorem C-1(i), (iv) of \fn{C1_identification_v44.tex} \tH\ (the step-4 solution is the smooth Lean flow from $\Gamma_0$ on $[0,t(27.15)]\ni t_1$; this replaces step~4's sentence ``smooth for $t>0$ by parabolic regularity'', \fn{step4_proof.tex} l.672, labelled \textup{[A]} there).
\item \emph{Symmetry.} $\gamma(t_1)$ is symmetric about the axis $\mathbb Re_1$ through the tip: $\sigma\circ\gamma(t_1)=\gamma(t_1)\circ\iota$ for the reflection $\sigma(z)=\bar z$ (axis taken as $\mathbb R$) and an involutive diffeomorphism $\iota$ of $\mathbb S^1$. (Theorem C-1(i) of \fn{C1_identification_v44.tex}, with $\iota(z)=\bar z$ on $\mathbb S^1\subset\mathbb C$ \tH.)
\item \emph{(H1)} The modulated decomposition exists at $\tau_1$: parameters $P_1=(t_1,\lambda_1,\xi_1,\alpha(\tau_1))$ with the three constraints satisfied, $|\alpha(\tau_1)|\le\alpha_1$, $\sqrt{\Lambda(\tau_1)}\le5\cdot10^{-8}$. This includes the chart statement (GS) of Definition~\ref{U:def:charts} at $\tau_1$, certified on $|x|\le60\frac14$ by Lemma GS-1 of \fn{GS_lemma_v44.tex} (\fn{fix14_interface/GS_tau1.json}: $\sup|ku|\le6.52\cdot10^{-3}$, $\sup|u'|\le0.0542$) \tI; and (G) at $\tau_1$ (graph over the ray on $y\in[14,700]$; \fn{continuation_v33.tex} Lemma~G and its gap item~2) \tI.
\end{enumerate}

\subsection{From the Lean development (existence, uniqueness, maximal flow, blow-up criterion)}\label{U:sec:lean}
\emph{Source} (text v42; independent review \fn{writeup/lean_bridge/REVIEW_W7.md}: ``sound; no blocking finding'', for the theorems of \fn{writeup/lean_bridge/snap_final/}; text v43: the S3 theorems (PC), (MS), (GE) of \fn{writeup/lean_bridge/snap_s3/}, written after review W7, were reviewed in \fn{fix23_delta_check/DELTA_CHECK_v42.md} Part~A: statements faithful, not vacuous). Micah Warren's Lean development \fn{CWFlow} (local repository \fn{agents/lean/CWFlow}; library at HEAD \fn{8718f95}, its dependency \fn{MM} at \fn{7b15b01}, toolchain \fn{leanprover/lean4:v4.33.1}), new directory \fn{CWFlow/Curve/} (16 files, 23 Nov 2026: the 14 files of the snapshot \fn{writeup/lean_bridge/snap_final/}, of which \fn{All.lean} is superseded, plus \fn{Periodic2Circle.lean}, \fn{MaxShift.lean} and the new aggregator \fn{All.lean} of the snapshot \fn{writeup/lean_bridge/snap_s3/}, written after review W7 for its items S3a, S3b). \textbf{Commit pending:} \fn{CWFlow/Curve/} is not yet committed (and the repository has no remote), so its files are identified by the SHA256 sums in \fn{writeup/lean_bridge/snap_final/SHA256SUMS} and \fn{writeup/lean_bridge/snap_s3/SHA256SUMS}; every cited file of the snapshot is byte-identical to the repository file (re-checked by \fn{tools/check_numbers_v44.py} of this text set). The theorems below are in \fn{Criterion.lean} (\fn{fe567708...}), \fn{N1.lean} (\fn{f1d3d641...}), \fn{N1Basic.lean} (\fn{357cf86a...}), \fn{Lift.lean} (\fn{0833e4df...}), \fn{Lift3.lean} (\fn{95d693ca...}), \fn{Energy.lean} (\fn{dbd9441f...}), \fn{Defs.lean} (\fn{e7598375...}), \fn{Periodic2Circle.lean} (\fn{f7d02857...}) and \fn{MaxShift.lean} (\fn{80b86094...}); aggregator \fn{All.lean} (\fn{91d5788e...}). Text v41 cited commit \fn{d72b744} of \fn{HSLflow-Whitney}; that commit does not contain these files, and it is no longer cited. \emph{Compile status:} each file of \fn{CWFlow/Curve/} (the aggregator \fn{All.lean} too, after the S3 additions) was compiled with \fn{lean}, one file at a time in dependency order, against the repository's existing build outputs (\fn{.olean}) of the imported library modules (\fn{writeup/lean_bridge/IMPL_LOG.md}; raw logs with return code 0 for \fn{LiftV}, \fn{Lift3}, \fn{N1}, \fn{Lift}, \fn{Criterion}; no log is kept for \fn{N1Basic}, whose build output exists and against which these compiled; text v43 (delta check fix23 S-1): the compile records of the S3 files \fn{Periodic2Circle}, \fn{MaxShift} and \fn{All} are the raw logs \fn{writeup/lean_bridge/raw_Periodic2Circle.log}, \fn{writeup/lean_bridge/raw_MaxShift.log} and \fn{writeup/lean_bridge/raw_All.log}; text v42 had no raw log for them). A full \fn{lake build CWFlow.Curve.All} is \emph{pending} (review W7 item S1); its log is to be kept (records item L-2 of \fn{RECORDS_ERRATA_v44.md}). Each theorem (L1)--(Sh), (PC), (MS), (GE) below is followed, in its file or (for (T4), (T5)) in \fn{Criterion.lean}, by an axiom check (\fn{guard_msgs} around \fn{print axioms}) whose expected output is \fn{[propext, Classical.choice, Quot.sound]}, so the compile fails if any other axiom, in particular \fn{sorryAx}, is used anywhere in the dependencies of these theorems (text v44, check fix25 F-3: \fn{print axioms} covers the dependencies of the named theorem, not every imported declaration); (Gl) is used in the proof of (T2) and is covered by its check. No \fn{sorry}, \fn{axiom} or \fn{native_decide} occurs in \fn{CWFlow/Curve/} (text search, review W7 \S2) \tR.

\emph{Setting} (\fn{Defs.lean}, \fn{N1Basic.lean}, \fn{N1.lean}). $M=\mathbb S^1$ is Mathlib's \fn{Circle}, the unit complex numbers, a smooth manifold modelled on \fn{E1 := EuclideanSpace R (Fin 1)} (text v41 said ``the unit sphere of \fn{EuclideanSpace R (Fin 2)}''; that was wrong); $V=\mathbb C$; $J=\fn{Jc}$, multiplication by $i$ (\fn{Jc_isom}: an isometry; \fn{Jc_JJ}: $J^2=-1$). On \fn{Circle}, $\iota(p)=p^{-1}=\bar p$. For $F:\mathbb S^1\to\mathbb C$, $\fn{lift}\,F(\theta)=F(e^{i\theta})$ (\fn{Circle.exp}), a $2\pi$-periodic map. For $c:\mathbb R\to\mathbb C$: $\fn{curv}\,c=\ip{c''}{ic'}/|c'|^3$ (the signed curvature, with $N=iT$), $\fn{dS}\,c\,f=f'/|c'|$ (the derivative $\partial_s$), and
\[
\fn{bendingEnergy}\,c=\int_0^{2\pi}(\fn{curv}\,c)^2\,|c'|\,d\theta ,
\]
so that $\fn{bendingEnergy}(\fn{lift}F)=\int k^2ds$ over the closed curve $F$ traversed once. \fn{IsLagImmersion E1 Jc F}: $F$ is $C^\infty$ in Mathlib's charts of \fn{Circle}, with injective derivative; at $n=1$ the Lagrangian clause is automatic, and \fn{isLagImmersion_iff} identifies the notion with \fn{IsCurveImmersion} (chartwise $C^\infty$ with injective derivative). There is no embeddedness, orientation or turning-number hypothesis. \fn{IsFlow E1 Jc T Phi}, for $\Phi:\mathbb R\to\mathbb S^1\to\mathbb C$: $(t,p)\mapsto\Phi(t,p)$ is $C^\infty$ on $[0,T]\times\mathbb S^1$, each $\Phi(t)$, $t\in[0,T]$, is an \fn{IsLagImmersion}, and in every chart $\partial_t\Phi=\fn{flowVel}\,J\,(\Phi(t))$ (one-sided at $t=0$ and $t=T$), where, for $X:E\to V$ with $g=X^*\ip\cdot\cdot$,
\[
H=\mathrm{tr}_g(D^2X)^\perp,\qquad \fn{flowPot}=\mathrm{tr}_g\ip{D(JH)}{DX},\qquad\fn{flowVel}=J\,\nabla_g\fn{flowPot}.
\]
$\Phi$ is a total function of $t\in\mathbb R$; its values outside the interval named are arbitrary and never used.
\begin{enumerate}[label=(L\arabic*),leftmargin=3em]
\item \fn{CWFlow.Curve.cdf_exists} (\fn{N1.lean}): if \fn{IsCurveImmersion F}, there are $T>0$ and $\Phi$ with $\Phi(0)=F$ and \fn{IsFlow E1 Jc T Phi} \tL\ (the library theorem \fn{exists_flow_of_immersion} at $n=1$).
\item \fn{CWFlow.Curve.cdf_unique} (\fn{N1.lean}): if \fn{IsFlow E1 Jc T Phi1}, \fn{IsFlow E1 Jc T Phi2} and $\Phi_1(0)=\Phi_2(0)$, then $\Phi_1(t)=\Phi_2(t)$ for all $t\in[0,T]$ \tL\ (equality of maps; the library theorem \fn{flow_unique_of_immersion} at $n=1$). For flows with a tangential term, \fn{cdf_unique_up_to_reparam} gives equality up to a family of diffeomorphisms $\chi_t$ with $\chi_0=\mathrm{id}$ \tL.
\item[(T1)] \fn{CWFlow.Curve.cdf_exists_maximal} (\fn{Criterion.lean}). If \fn{IsLagImmersion E1 Jc F0}, there are $T_{\max}\in(0,\infty]$ and $\Phi$ with $\Phi(0)=F_0$ such that: (i) \fn{IsFlow E1 Jc T Phi} for every real $T$ with $0<T<T_{\max}$; (ii) for every $T>0$ and every $\Psi$ with \fn{IsFlow E1 Jc T Psi} and $\Psi(0)=F_0$: $T<T_{\max}$ and $\Psi(t)=\Phi(t)$ for $t\in[0,T]$; (iii) if $T_{\max}<\infty$, then for every $K\in\mathbb R$ there is $t\in[0,T_{\max})$ with $\fn{bendingEnergy}(\fn{lift}(\Phi(t)))>K$ \tL. ($T_{\max}$ is an extended nonnegative real; its real value is used only under $T_{\max}<\infty$.)
\item[(T2)] \fn{CWFlow.Curve.cdf_extends_of_bounded_energy} (\fn{Criterion.lean}). If $0<T$, \fn{IsFlow E1 Jc T' Phi} for every $T'\in(0,T)$, $t_0\in(0,T)$, and $\fn{bendingEnergy}(\fn{lift}(\Phi(t)))\le K$ for $t\in[t_0,T)$, then there are $\varepsilon>0$ and $\Psi$ such that $\Psi$ is an \fn{IsFlow E1 Jc} on $[0,T+\varepsilon]$ and $\Psi(t)=\Phi(t)$ for $t\in[0,T)$ \tL. (T1) is proved from (T2) in Lean.
\item[(T3)] \fn{CWFlow.Curve.aNormSq_eq_curv_sq} (\fn{Lift3.lean}): in a chart of \fn{Circle}, the library's $|A|^2$ of an immersion $F$ equals $\fn{curv}(\fn{lift}F)(\theta)^2$ at the corresponding $\theta$ \tL. Used only inside the Lean proof of (T2); not an input of this text.
\item[(T4)] \fn{CWFlow.Curve.IsFlow.lift_isCDF} (\fn{Lift.lean}): if \fn{IsFlow E1 Jc T Phi}, then $\gamma(t,\theta)=\Phi(t,e^{i\theta})$ is \fn{IsCDF} on $(0,T)$: $C^\infty$ on $(0,T)\times\mathbb R$, $2\pi$-periodic in $\theta$, $\gamma_\theta\ne0$, and $\partial_t\gamma=-(\partial_s\partial_s\kappa)\,i\gamma_\theta/|\gamma_\theta|$ with $\kappa=\fn{curv}$, $\partial_s=\fn{dS}$ \tL. (Interior times only; at $t=0$ and $t=T$ the velocity identity is the clause of \fn{IsFlow} together with (V).)
\item[(T5)] \fn{CWFlow.Curve.IsCDF.curv_sq_le} (\fn{Energy.lean}): an \fn{IsCDF} on $(a,b)$ with $\fn{bendingEnergy}\le K$ on $[t_0,b)$ has $\kappa^2\le C$ on $[t_0,b)\times\mathbb R$ for some $C$ \tL. Used only inside the Lean proof of (T2).
\item[(V)] \fn{CWFlow.Curve.flowVel_eq_neg_kss} (\fn{N1Basic.lean}): for $X:\mathbb R\to\mathbb C$ of class $C^\infty$ with $X'\ne0$ everywhere, $\fn{flowVel}\,\fn{Jc}\,X=-(\partial_s\partial_s\kappa)\,iX'/|X'|$ with $\kappa=\fn{curv}X$, $\partial_s=\fn{dS}X$ \tL.
\item[(Sy)] \fn{CWFlow.Curve.cdf_symmetric} (\fn{N1.lean}): if \fn{IsFlow E1 Jc T Phi} and $\overline{\Phi(0,p^{-1})}=\Phi(0,p)$ for all $p\in\mathbb S^1$, then $\overline{\Phi(t,p^{-1})}=\Phi(t,p)$ for all $t\in[0,T]$ and $p\in\mathbb S^1$ \tL\ (from \fn{IsFlow.conj_reflect} and (L2)).
\item[(Sh)] \fn{CWFlow.Curve.IsFlow.shift_gen} (\fn{N1.lean}): if \fn{IsFlow E J T Phi} and $a\ge0$, then $t\mapsto\Phi(t+a)$ is an \fn{IsFlow} on $[0,T-a]$ \tL.
\item[(Gl)] \fn{CWFlow.rst_glue} (library, \fn{Flow/Restart.lean}): if $a\ge0$, $0<s\le T_1$, $\Phi$ is an \fn{IsFlow} on $[0,a+s]$, $\Theta$ is an \fn{IsFlow} on $[0,T_1]$ and $\Theta(t)=\Phi(a+t)$ for $t\in[0,s]$, then there is an \fn{IsFlow} $\Psi$ on $[0,a+T_1]$ with $\Psi=\Phi$ on $[0,a+s]$ \tL\ (gluing with an overlap of positive length).
\item[(PC)] \fn{CWFlow.Curve.exists_isLagImmersion_of_periodic} (\fn{Periodic2Circle.lean}): if $c:\mathbb R\to\mathbb C$ is $C^\infty$, $2\pi$-periodic and $c'(\theta)\ne0$ for all $\theta$, there is $F:\mathbb S^1\to\mathbb C$ with $\fn{lift}\,F=c$ and \fn{IsLagImmersion E1 Jc F}; and \fn{lift_injective}: $F$ is determined by $\fn{lift}\,F$ \LeanSthree\ (review W7 S3a; applied to one regular periodic curve, e.g.\ $c(\theta)=e^{i\theta}$, with $c'(\theta)=ie^{i\theta}\ne0$, it shows that the hypothesis of (T1) is inhabited).
\item[(MS)] \fn{CWFlow.Curve.cdf_maximal_shift} (\fn{MaxShift.lean}): let $\Phi$, $T_{\max}$ satisfy clauses (i), (ii) of (T1) for $F_0=\Phi(0)$, and let $0\le t_1$, $t_1<T_{\max}$. Then $T_{\max}-t_1>0$; $t\mapsto\Phi(t_1+t)$ is an \fn{IsFlow} on every $[0,T]$ with $0<T<T_{\max}-t_1$; every \fn{IsFlow} $\Psi$ on $[0,T]$, $T>0$, from $\Phi(t_1)$ has $T<T_{\max}-t_1$ and equals $\Phi(t_1+\cdot)$ there; and if $T_{\max}-t_1<\infty$, then $\fn{bendingEnergy}(\fn{lift}(\Phi(t_1+t)))$ is unbounded for $t\in[0,T_{\max}-t_1)$ \LeanSthree\ (review W7 S3b).
\item[(GE)] \fn{CWFlow.Curve.isFlow_glue_endpoint} (\fn{MaxShift.lean}): if $a,b>0$, $\Phi$ is an \fn{IsFlow} on $[0,a]$, $\Theta$ is an \fn{IsFlow} on $[0,b]$ and $\Theta(0)=\Phi(a)$, then there is an \fn{IsFlow} $\Psi$ on $[0,a+b]$ with $\Psi=\Phi$ on $[0,a]$ and $\Psi(a+t)=\Theta(t)$ for $t\in[0,b]$ \LeanSthree\ (review W7 S3b, endpoint gluing).
\end{enumerate}
\emph{Hypotheses carried, junk values.} None of these theorems carries an unproved hypothesis: the library's hypotheses (\fn{hJi}, \fn{hJJ}, \fn{hdim}) are discharged by \fn{Jc_isom}, \fn{Jc_JJ}, \fn{Complex.finrank_real_complex}, and there is no hypothesis $\dim E\ge2$ \tR\ (review W7 \S1). Every \fn{curv} and \fn{bendingEnergy} in (T1)--(T5) is taken at a time at which the curve is an immersion, where \fn{lift} is $C^\infty$ with nonvanishing derivative (\fn{contDiff_lift_of_chart}, \fn{deriv_lift_ne_zero_of_chart}), so Mathlib's junk values of \fn{deriv} and of the interval integral never enter \tR\ (review W7 \S2). $\int\kappa^2|c'|\,d\theta$ is invariant under reparametrisation, so $\fn{bendingEnergy}(\fn{lift}(\Phi(t)))=\int_{\gamma(t)}k^2ds$ for every parametrisation of the same closed curve \tH.

\begin{lemma}[$n=1$: the Lean velocity is CDF]\label{U:lem:vel}
Let $X$ be (a) a $C^\infty$ map $X:\mathbb R\to\mathbb C$ with $X_\theta\ne0$ on all of $\mathbb R$, or (b) the lift (\fn{lift}) of an immersion of \fn{Circle}, read in a chart (text v44, check fix25 F-5: these are the two cases that (V) and \fn{flowVel_chart_eq_lift} cover, and nothing else is used). Then $\fn{flowVel}\,J\,X=-k_{ss}N$, with $T=X_\theta/|X_\theta|$, $N=JT$, $k=\ip{X_{\theta\theta}}{JX_\theta}/|X_\theta|^3$, $\partial_s=|X_\theta|^{-1}\partial_\theta$. Moreover $\fn{flowVel}(\sigma X)=\sigma\,\fn{flowVel}(X)$ for $\sigma(z)=\bar z$, and $\fn{flowVel}(X\circ\varphi)=\fn{flowVel}(X)\circ\varphi$ for every diffeomorphism $\varphi$ of $\mathbb R$ (either orientation).
\end{lemma}
\begin{proof}\tL\ (text v42). (V) \fn{flowVel_eq_neg_kss}; for a chart representative of an immersion of \fn{Circle}, the velocity is that of the lift (\fn{flowVel_chart_eq_lift}, \fn{LiftV.lean}; for a flow, (T4)); $\sigma$: \fn{flowVel_conj_any} (\fn{N1Basic.lean}, no hypotheses); reparametrisation: the library's \fn{flowVel_comp} (an invertible derivative suffices, so either orientation). The exact sympy identities S1--S3 of \fn{U1_sym.py} \tS\ remain as an independent cross-check.
\end{proof}
\noindent So an \fn{IsFlow} at $n=1$ is a smooth closed immersed curve moving with velocity $-k_{ss}N$ and \emph{no tangential part}: the parametrised curve diffusion flow $\partial_tf=-\nabla_s^2\kappa$, i.e.\ normal velocity $-\kappa_{ss}$ as in step~3.

\subsection{Dziuk--Kuwert--Sch\"atzle 2002 (no longer an input)}
Text v42: (D1) and (D2) of text v41 are not used. The criterion is (T1)(iii) and the extension form is (T2) \tL, so the author item A-2 (that ``maximal solution'' in the sense of DKS means maximal among smooth solutions) and the re-reading of DKS for (D1) are void. As literature only: the $L^2$ criterion for curve diffusion is Theorem~3.1 of Dziuk, Kuwert and Sch\"atzle (2002), which also gives the rate $\int\kappa^2ds\ge c(T-t)^{-1/4}$ (not used: (T1) gives $\sup=\infty$ without a rate, and every use in the step-3 texts is in the sup form), and, for plane curves of nonzero area, Chou (2003). \fn{writeup/lit_P1/P1_CONFIRMATION.md} is now a historical record. No embeddedness, area or turning-number hypothesis is involved anywhere.

\subsection{From step 3 (objects; all \tR)}
$S_1$ solves a polynomial ODE, so it is real-analytic. $\hat\psi_3,\hat h_2$ are the $C^4$ functions of \fn{step3nl_v4c/agent_int/v3/c4fun.py}: finitely many polynomial pieces with degree-9 Hermite corrections matching jets $0..4$ at every interface, and an explicit smooth extension for $x\ge28$; so they are $C^4$, $C^\infty$ off a finite set $B_0\subset[0,28]$, with bounded one-sided fifth derivatives, even. The cutoffs: $\chi=\chi_{32}$ and $\zeta$ (rising on $[10.25,16.25]$, falling on $[57,60]$) are built from $S_7(t)=t^4(35-84t+70t^2-20t^3)$, so they are $C^3$ with $\chi^{(j)},\zeta^{(j)}$ ($j=1,2,3$) vanishing at their four break points and $|\chi''''|,|\zeta''''|$ bounded with jumps there ($S_7''''(0)=840$, $S_7''''(1)=-840$ \tS\ S6); (A)'s cutoffs $\eta_{\rm in},\eta_o$ use $S(t)=t^3(10-15t+6t^2)$: $C^2$, with $S'''$ jumping by $60$ at both ends \tS\ S6. Weights: $\rho=e^{-\phi}$ is smooth; $W=e^{-\psi}$ with $\psi'=(1-s)\frac{4x}{1+x^2}+s\phi'$, $s=S_7((x-30.25)/8)$, so $\psi\in C^4$ and $W\in C^4$; $\varepsilon\omega=28y^{-5}$ is smooth on $y>0$. The graph equation for the normal graph $u$ over $S_1$ in the moving frame (\fn{step3nl_v4c/BOOTSTRAP.md} ``Unknowns and equations'', \fn{graph_eq.py}, \fn{graph_struct.pkl}) is taken as derived there (exact sympy, not re-derived here) \tR; its frame part is the identity (S5) below. The modulation system and the invertibility of its matrix on the validity box are those of \fn{agent_mod/ibp/MOD_IBP.md} (Neumann step $\nu\le0.0466<1$ under $|\alpha|\le5\cdot10^{-4}$, $A,W_1\le10^{-3}$, strip $|w|,|w'|\le0.02$) \tR.

%======================================================================
\section{The physical solution}\label{U:sec:phys}
\begin{lemma}[gluing]\label{U:lem:glue}
If $a,b>0$, $\Phi_1$ is an \fn{IsFlow} on $[0,a]$, $\Phi_2$ is an \fn{IsFlow} on $[0,b]$ and $\Phi_2(0)=\Phi_1(a)$, then there is an \fn{IsFlow} $\Psi$ on $[0,a+b]$ with $\Psi=\Phi_1$ on $[0,a]$ and $\Psi(a+t)=\Phi_2(t)$ for $t\in[0,b]$. In particular the piecewise map $\Phi(t)=\Phi_1(t)$ ($t\le a$), $\Phi_2(t-a)$ ($t\ge a$) equals $\Psi$ on $[0,a+b]$ and is an \fn{IsFlow} there.
\end{lemma}
\begin{proof}\LeanSthree\ The first sentence is (GE) \fn{CWFlow.Curve.isFlow_glue_endpoint}. The second is a hand step \tH: $\Psi=\Phi$ pointwise on $[0,a+b]$, and the clauses of \fn{IsFlow} on $[0,a+b]$ (joint smoothness on $[0,a+b]\times\mathbb S^1$, immersion at each time, the equation with time derivatives within $[0,a+b]$) involve only the values on $[0,a+b]$ (text v43, delta check fix23 W-4: text v42 labelled this congruence step [Lean] too). The overlap form is (Gl) \fn{rst_glue}. (Text v42: text v41 had a hand proof. Lemma~\ref{U:lem:max} no longer needs gluing: it takes maximality from (T1).)
\end{proof}

\begin{lemma}[maximal smooth solution]\label{U:lem:max}
Let $F=\gamma(t_1)$. There is $T_{\max}\in(0,\infty]$ and a unique $\Phi:[0,T_{\max})\times\mathbb S^1\to\mathbb C$ such that $\Phi|_{[0,T]}$ is an \fn{IsFlow} with $\Phi(0)=F$ for every $T<T_{\max}$, and there is no \fn{IsFlow} on $[0,T_{\max}]$ with $\Phi(0)=F$ when $T_{\max}<\infty$. Every \fn{IsFlow} on $[0,T]$ from $F$ is the restriction of $\Phi$. Put $\gamma(t_1+t):=\Phi(t)$.
\end{lemma}
\begin{proof}\tL\ (text v42). $F=\gamma(t_1)$ is the value at $t_1$ of an \fn{IsFlow} on an interval $[0,T]\ni t_1$ (Theorem C-1 (i), (iv) of \fn{C1_identification_v44.tex}), so it is an \fn{IsLagImmersion E1 Jc} by the immersion clause of \fn{IsFlow}; no conversion from a periodic curve is needed here. (T1) with $F_0=F$ gives $T_{\max}$ and $\Phi$. Clause (i) is ``$\Phi|_{[0,T]}$ is an \fn{IsFlow} for every $T<T_{\max}$''. Clause (ii) gives that every \fn{IsFlow} from $F$ on $[0,T]$ has $T<T_{\max}$ and is the restriction of $\Phi$; in particular there is none on $[0,T_{\max}]$ when $T_{\max}<\infty$, and $\Phi$ is unique on $[0,T_{\max})$. (Text v41 assembled this by hand from (L1), (L2) and the gluing lemma.)
\end{proof}

\begin{lemma}[symmetry]\label{U:lem:sym}
$\sigma\circ\gamma(t)=\gamma(t)\circ\iota$ for all $t\in[t_1,t_1+T_{\max})$.
\end{lemma}
\begin{proof}\tL\ (text v42). (Sy) \fn{cdf_symmetric}, applied to $\Phi|_{[0,T]}$ for each $T<T_{\max}$; its hypothesis $\overline{\Phi(0,p^{-1})}=\Phi(0,p)$ is (S4.2) with $\iota(p)=p^{-1}=\bar p$ on \fn{Circle}, and its conclusion $\overline{\Phi(t,p^{-1})}=\Phi(t,p)$ for all $p$ is the claim (replace $p$ by $p^{-1}$). (Text v41 had a hand proof through S2, S3; the closure of \fn{IsFlow} under $\sigma\circ\cdot\circ\iota$ is now the Lean lemma \fn{IsFlow.conj_reflect}.)
\end{proof}

\begin{lemma}[blow-up criterion, the form needed]\label{U:lem:bu}
If $T_{\max}<\infty$ then $\sup_{t<t_1+T_{\max}}\int_{\gamma(t)}k^2ds=\infty$.
\end{lemma}
\begin{proof}\tL\ (text v42). Clause (iii) of (T1): if $T_{\max}<\infty$, then $\fn{bendingEnergy}(\fn{lift}(\Phi(t)))$ is unbounded on $[0,T_{\max})$, and this quantity is $\int_{\gamma(t_1+t)}k^2ds$ (\S\ref{U:sec:lean}, last paragraph before Lemma~\ref{U:lem:vel}). (Texts v40, v41 derived this from the statement of DKS Theorem~3.1, with the author item A-2 on ``maximal solution''; neither is used now.)
\end{proof}
\begin{remark}[extension form; (D1) replaced]\label{U:rem:D1}
The extension form of the criterion is (T2) \fn{cdf_extends_of_bounded_energy} \tL: if $\int k^2ds$ stays bounded near a finite end, the flow extends to an \fn{IsFlow} on a closed interval beyond it, so the extension is an immersion at the end time and satisfies the velocity identity there (the ``C-4 completion'' that text v41 added to DKS (D1) by hand is part of the Lean statement). (T1) is proved from (T2) in Lean; nothing in U-1 uses (T2) directly.
\end{remark}
\noindent This is (P1$'$), with existence, uniqueness, maximality, symmetry and the criterion from the compiled Lean theorems (T1), (Sy) (and (L1), (L2)) \tL, subject to the build and commit items of \S\ref{U:sec:gaps} (text v42; text v41: (L1), (L2) read from source and DKS (D2)).

%======================================================================
\section{Charts, frame and modulation}\label{U:sec:mod}
\begin{definition}[frame and charts]\label{U:def:charts}
For $P=(t,\lambda,\xi,\alpha)$ with $t\in[t_1,t_1+T_{\max})$, $\lambda>0$, $\xi\in\mathbb R$ put $\Gamma_P:=(\gamma(t)-\xi e_1)/\lambda$. Let $\Psi(x,n)=S_1(x)+n\nu_{S_1}(x)$ ($x$ arclength from the tip); $D\Psi$ is invertible where $1-k(x)n\ne0$.
\begin{itemize}[leftmargin=*]
\item \emph{(GS) at $P$}: there is $u\in C^1([-X_S,X_S])$, $X_S:=60$, with $\sup|k\,u|\le\frac12$, $\sup|u'|\le1$, such that $x\mapsto\Psi(x,u(x))$ reparametrises a sub-arc of $\Gamma_P$ containing the tip (the arc $\Gamma_{\rm tip}$ of \fn{continuation_v44.tex} \S2). \emph{(GS$^\circ$) at $P$}: the same with $\sup|ku|<1$, $\sup|u'|<2$ (open, relaxed; \fn{GS_lemma_v44.tex} Definition GS). Text v37 had $X_S=60\frac14$; nothing in this file uses $u$ on $(60,60\frac14]$ (\fn{GS_lemma_v44.tex} \S1; review fix15 D-1).
\item \emph{(G) at $P$}: the ray-graph statement of \fn{continuation_v33.tex} Definition~G on $y\in[14,700]$ (slope $|f|\le\frac12$).
\end{itemize}
\end{definition}

\begin{lemma}[smooth dependence of the graph functions]\label{U:lem:graph}
The set $\mathcal O_{GS}$ of $P$ at which (GS$^\circ$) holds is open, and on it $u=\mathcal U(t,\lambda,\xi;x)$ is $C^\infty$ jointly in $(t,\lambda,\xi,x)$; the same holds for (GS) with a gap (Lemma GS-3 of \fn{GS_lemma_v44.tex}, which treats the endpoints $x=\pm60$). The same holds for the ray-graph height $h=\mathcal H(t,\lambda,\xi;y)$ and its slope $f=\partial_y\mathcal H$ under (G), and both are even in the sense $\mathcal U(\cdot;-x)=\mathcal U(\cdot;x)$.
\end{lemma}
\begin{proof}\tH\ $G(t,\lambda,\xi,\theta):=(\Phi(t-t_1,\theta)-\xi e_1)/\lambda$ is $C^\infty$ on $[t_1,t_1+T)\times(0,\infty)\times\mathbb R\times\mathbb S^1$ for every $T<T_{\max}$ (Lemma~\ref{U:lem:max}). Write $\Psi^{-1}\circ G=(X,N)$ near the arc (well defined and smooth where $|kN|<1$, locally). Under (GS), $\theta\mapsto X(P,\theta)$ is a $C^1$ bijection from the parameter arc onto $[-X_S,X_S]$ with $C^1$ inverse $\theta(x)$; differentiating $X(P,\theta(x))=x$ gives $\partial_\theta X\ne0$. The implicit function theorem applied to $X(P,\theta)-x=0$ gives $\theta=\Theta(P,x)$ smooth near each point, hence $\mathcal U=N(P,\Theta(P,x))$ smooth; the local solutions agree on overlaps by uniqueness of $\theta(x)$, and compactness of $[-X_S,X_S]$ gives a uniform neighbourhood of $P$ on which (GS) holds strictly. Evenness: Lemma~\ref{U:lem:sym}, $\xi e_1$ on the axis, and $S_1$ even. The ray chart is the same argument with $\Psi$ replaced by the affine normal coordinates of the ray.
\end{proof}

\noindent\emph{Frame and modulation equations} (convention of \fn{step3_proof_v32.tex} \S1.1). Along a $C^1$ path $P(\tau)$ with
\begin{equation}\label{U:eq:ode}
t'=\lambda^4,\qquad\lambda'=-(\tfrac14-a)\lambda,\qquad\xi'=-\lambda b,\qquad\alpha'=\tilde\lambda\alpha+\beta ,
\end{equation}
$\Gamma(\tau):=\Gamma_{P(\tau)}$ satisfies $\ip{\Gamma_\tau}{\nu}=-\kappa_{ss}+(\frac14-a)\ip\Gamma\nu+b\ip{e_1}\nu$ (\tS\ S5 with the scaling of $\kappa_{ss}$; the velocity of $\gamma$ is $-k_{ss}N$ by Lemma~\ref{U:lem:vel}). The sign of $b$ is fixed by this formula; the verifier uses only $|a|,|b|,|\beta|$. With $w(\tau,\cdot)=u-\alpha\hat\psi_3-\alpha^2\hat h_2$, $\tilde w=\chi w$, the constraints are $c_j(P):=\ip{\chi(\mathcal U(P;\cdot)-\alpha\hat\psi_3-\alpha^2\hat h_2)}{f_j}_\rho=0$, $j=1,2,3$, with the fixed functions $f_j$ of \fn{step3_proof_v32.tex} \S3.1.

\begin{proposition}[(P2$'$): modulation]\label{U:prop:mod}
Let $\mathfrak B$ be the closed set $\mathfrak B_{\rm mod}$ of \fn{bootstrap_v44.tex} (Definition BS): the $P$ at which (GS) holds, the modulation validity box holds ($|\alpha|\le5\cdot10^{-4}$, $A,W_1\le10^{-3}$, $|w|,|w'|,|w''|\le0.02$ on $\sigma\in[29,32]$), and B1 holds in the form of the flux lemma, $\Lambda\le L(\tau)+2C'$ on the march and $\le\ell(\tau)+2C'$ in the tail (review fix15 D-5; text v37 had $\sqrt\Lambda\le\sqrt{\Lambda_{\rm bar}}$). Assume (S4.1)--(S4.3). Then:
\begin{enumerate}[label=(\alph*),leftmargin=2em]
\item on the open set $\mathcal O:=\{P\in\mathcal O_{GS}:\ (m_{ij}(P))\text{ invertible}\}\supset\mathfrak B$ (defined through (GS$^\circ$), review C-5) there are $C^\infty$ functions $(a,b,\beta)=\mathcal A(P)$ such that every $C^1$ solution of \eqref{U:eq:ode} with $(a,b,\beta)=\mathcal A(P)$ keeps $c_j(P(\tau))$ constant;
\item \eqref{U:eq:ode} with $(a,b,\beta)=\mathcal A(P)$ and $P(\tau_1)=P_1$ has a unique maximal $C^1$ solution $P:[\tau_1,\tau_{\rm mod})\to\mathcal O$, and $c_j\equiv0$ along it;
\item \emph{(continuation)} if $\tau_{\rm mod}<\infty$ and $P(\tau)\in\mathfrak B$ for all $\tau\in[\tau_1,\tau_{\rm mod})$, then $t(\tau)\uparrow t_1+T_{\max}$ as $\tau\uparrow\tau_{\rm mod}$; in that case $T_{\max}=\int_{\tau_1}^{\tau_{\rm mod}}\lambda^4d\tau$.
\end{enumerate}
So the modulation exists as long as the physical solution exists and the bootstrap box holds, and the end $\tau_{\max}$ of the rescaled interval in \fn{continuation_v33.tex} is $\tau_{\rm mod}$, with $t(\tau_{\max}^-)=t_1+T_{\max}$ whenever $\tau_{\max}<\infty$ and the bootstrap holds up to $\tau_{\max}$.
\end{proposition}
\begin{proof}\tH\ (a) By Lemma~\ref{U:lem:graph}, $c_j$ is $C^\infty$ in $(t,\lambda,\xi)$ and polynomial in $\alpha$ (the integrand is supported in $|x|\le34.25$, a compact set, and $f_j\rho$ is fixed). Along \eqref{U:eq:ode},
\[
\tfrac{d}{d\tau}c_j=r_j(P)+a\,m_{1j}(P)+b\,m_{2j}(P)+\beta\,m_{3j}(P),
\]
$r_j=\lambda^4\partial_tc_j-\frac\lambda4\partial_\lambda c_j+\tilde\lambda\alpha\,\partial_\alpha c_j$, $m_{1j}=\lambda\partial_\lambda c_j$, $m_{2j}=-\lambda\partial_\xi c_j$, $m_{3j}=\partial_\alpha c_j$, all $C^\infty$. Written with the graph equation \tR, $\sum_ix_im_{ij}=-r_j$ is the system $x^T(G+E(u))=R$ of \fn{MOD_IBP.md} / \fn{E0_LEMMA_v4c_v44.tex} (``Setting''), with $G$ the fixed Gram matrix and $E$ the state part; the Neumann bound $\|G^{-1}E\|\le\nu\le0.0466<1$ on the validity box \tR\ makes $(m_{ij})$ invertible on $\mathfrak B$; $\mathcal O$ is open because $\mathcal O_{GS}$ is open (Lemma~\ref{U:lem:graph}) and invertibility is an open condition on the continuous matrix $(m_{ij})$. Put $\mathcal A=-(m^T)^{-1}r$.
(b) The right side of \eqref{U:eq:ode} is $C^1$ on $\mathcal O$: Picard--Lindel\"of and the usual maximal extension. $c_j(P(\tau_1))=0$ by (S4.3) and $\frac{d}{d\tau}c_j=0$.
(c) Suppose not: $t$ is increasing ($t'=\lambda^4>0$), so $t(\tau)\to t_*<t_1+T_{\max}$. On $\mathfrak B$ the modulation lemma bounds $|a|,|b|,|\beta|$ by a constant $\bar c$ (the box bounds $A,W_1$, and B1 bounds $B=\|\tilde w''\|_\rho$; \fn{MOD_IBP.md} final inequalities \tR). Hence on $[\tau_1,\tau_{\rm mod})$: $\lambda\in[\lambda_1e^{-(\frac14+\bar c)(\tau_{\rm mod}-\tau_1)},\lambda_1e^{\bar c(\tau_{\rm mod}-\tau_1)}]$, $|\xi-\xi_1|\le\bar c\,\lambda_1e^{\bar c(\tau_{\rm mod}-\tau_1)}(\tau_{\rm mod}-\tau_1)$, $|\alpha|\le5\cdot10^{-4}$. The set $K:=\mathfrak B\cap\{t\in[t_1,t_*],\ \lambda,\xi\text{ in these ranges}\}$ is bounded, and closed: a limit of $P_n\in K$ has $t\le t_*<t_1+T_{\max}$, so $\Gamma_{P_n}\to\Gamma_P$ in $C^\infty$; the graph functions converge in $C^1$ with the same bounds, so (GS) holds at $P$ (the limit argument of Lemma GS-4 of \fn{GS_lemma_v44.tex}), and the box inequalities are closed conditions in $P$ (they are continuous functions of $P$ by Lemma~\ref{U:lem:graph}). So $K$ is a compact subset of $\mathcal O$ that contains $P(\tau)$ for all $\tau<\tau_{\rm mod}<\infty$, which is impossible for a maximal solution of a $C^1$ ODE. The last identity is $t'=\lambda^4$.
\end{proof}

%======================================================================
\section{The theorem}\label{U:sec:thm}
Let $B$ be the finite set of $x$ where one of $\hat\psi_3^{(5)},\hat h_2^{(5)},\chi'''',\zeta''''$ jumps ($B\subset[-34.25,34.25]\cup\{\pm57,\pm60\}$, even). $\mathrm{IR}=[10.25,16.25]$, $\mathrm{OR}=[57,60]$, $g=w'$, $v=\zeta g$, $D_2v=(Wv'')''/W$, and $V_j,E,E_I,\Lambda,A,W_1,B,W_3,D_4$, $E^A_j$, $z^A_k$, $N_l,M_l,s_2,K^{\rm OR}_0,K^{\rm OR}_2,\varphi_2$ as in \fn{EI_SPEC.md} \S1, \fn{lemmaS_A_v33.tex} \S1 and \fn{flux2_lemma_v44.tex} \S\S1--3.

\begin{theorem}[U-1]\label{U:thm}
Assume (S4.1)--(S4.3). Let $\tau'<\tau_{\rm mod}$ and assume (GS) and (G) at $P(\tau)$ for $\tau\in[\tau_1,\tau']$. (Text v40, review fix18 F-1: no other member of the bootstrap set is needed, in particular neither B1(b) nor the validity box; text v39 assumed $P(\tau)\in\mathfrak B$. The proof below uses only that $P$ is the $C^1$ solution of (P2$'$)(b) on $[\tau_1,\tau']\subset[\tau_1,\tau_{\rm mod})$, (GS) for the graph function $u=\mathcal U(P;\cdot)$, and (G) for the far-field energies of (d). Text v41, delta check \fn{fix20_delta_check/DELTA_CHECK_v40.md} W-2: this concerns the regularity statements; the inequalities quoted in (c) and (d) keep the hypotheses of their own lemmas, and U-1 only upgrades ``for a.e.\ $\tau$'' to ``for every such $\tau$''.) Then:
\begin{enumerate}[label=(\alph*),leftmargin=2em]
\item \emph{(unknowns)} $t,\lambda,\xi,\alpha\in C^1$ and $a,b,\beta\in C^1$ on $[\tau_1,\tau']$. $u(\tau,x)=\mathcal U(P(\tau);x)$: every $\partial_x^ku$ and $\partial_\tau\partial_x^ku=\partial_x^k\partial_\tau u$ is continuous on $[\tau_1,\tau']\times[-X_S,X_S]$, and $u$ is even in $x$. $w$: $\partial_x^kw$ and $\partial_\tau\partial_x^kw$ ($k\le4$) are continuous; $\partial_x^5w$ is continuous on $[\tau_1,\tau']\times([-X_S,X_S]\setminus B)$ and bounded. Hence $g,v\in C^3$ in $x$ with $g'''$, $v'''$ Lipschitz in $x$ uniformly in $\tau$, and $g'''',v''''$ bounded and continuous off $B$; on $\mathrm{OR}\cup[28,60]$ ($\hat\psi_3,\hat h_2$ smooth there) $g$ is $C^\infty$ in $x$. The $w$-equation and the $v$-equation (T1)--(T12) of \fn{EI_SPEC.md} \S2 hold pointwise on $[\tau_1,\tau']\times([-X_S,X_S]\setminus B)$.
\item \emph{(verifier energies)} $A^2,B^2,E,V_0^2,V_2^2,E_I,\Lambda\in C^1([\tau_1,\tau'])$, and for \emph{every} $\tau$:
$\frac12\frac{d}{d\tau}A^2=\ip{\chi w_\tau}{\tilde w}_\rho$; $\frac12\frac{d}{d\tau}B^2=\ip{(\rho\tilde w'')''/\rho}{\chi w_\tau}_\rho$ (no boundary term); $\frac12\frac{d}{d\tau}V_0^2=\ip{v}{v_\tau}_W$; $\frac12\frac{d}{d\tau}V_2^2=\ip{D_2v}{v_\tau}_W$ with $D_2v=v''''-2\psi'v'''+(\psi'^2-\psi'')v''$ \tS\ S7 (defined off $B$, bounded); the level-0 integration by parts $\ip{v}{\zeta f'}_W=-\int(W\zeta v)'f$ for every $f\in C^1$ \tS\ S7; and the integrations by parts of \fn{MOD_IBP.md} (test functions $\Psi=\chi f_j\rho\in C^3$).
\item \emph{(flux lemma)} for every $\tau$: $g(\tau,\cdot)\in C^\infty(\mathrm{OR})$, so $g''''(\tau,\cdot)\in L^2(\mathrm{OR})$; $M_l,N_l$ ($l\le4$), $s_2$, $K^{\rm OR}_0,K^{\rm OR}_2$, $\varphi_2$ are continuous on $[\tau_1,\tau']$, hence $\varphi=\mu\varepsilon_I\varphi_2\in L^1$ of every window; Agmon $s_2^2\le\frac13M_2^2+2M_2M_3$ holds for every $\tau$; wherever the hypotheses of the $E_I$ and flux lemmas (B1(c), BF, B2$'$) hold, the inequality (F:eq:star) of \fn{flux2_lemma_v44.tex} holds at every such $\tau$, not only a.e.; and the hypotheses ``$\Lambda$ absolutely continuous, $\varphi\ge0$ integrable'' of the absorption lemma hold.
\item \emph{(Lemma S(A), Lemma U)} $E^A_j\in C^1([\tau_1,\tau'])$, $z^A_k$ is continuous in $(\tau,s)$ and $C^\infty$ in $s$ on $[\ln14,\ln700]$, the identities of \fn{lemmaS_A_v33.tex} \S2 (with the point terms of its Lemma~P) hold at every $\tau$, so, wherever the hypotheses of Lemma S(A) hold, the level-$j$ inequality (D1) there holds at every such $\tau$, not only a.e.; the window integrals $X_k(\tau_a,T)$ and $\int_{\tau_a}^{\tau_a+T}V_m^2$ are continuous in $(\tau_a,T)$. The same regularity holds for every energy of the form $\int\Theta(\tau,s)(\partial_s^kv)^2ds$ with $\Theta$ continuous, $C^1$ in $\tau$ and compactly supported in the region where (G) holds (the window energies of Lemma~U and the zone-III energies).
\end{enumerate}
\end{theorem}

\subsection{Two calculus facts}
\begin{lemma}[differentiation]\label{U:lem:leib}
Let $K$ be compact, $\Theta,\phi:[\tau_1,\tau']\times K\to\mathbb R$ continuous with $\partial_\tau\Theta,\partial_\tau\phi$ continuous on $[\tau_1,\tau']\times(K\setminus B')$, bounded, $B'$ finite. Then $F(\tau)=\int_K\Theta\phi^2$ is $C^1$ with $F'=\int_K(\partial_\tau\Theta\,\phi^2+2\Theta\phi\,\partial_\tau\phi)$.
\end{lemma}
\begin{proof}\tH\ $F(\tau+h)-F(\tau)=\int_K\int_0^1\partial_\tau(\Theta\phi^2)(\tau+\vartheta h,x)\,d\vartheta\,h\,dx$ (fundamental theorem of calculus in $\tau$ for each $x\notin B'$); divide by $h$ and use dominated convergence; the derivative is continuous in $\tau$ by dominated convergence again.
\end{proof}
\begin{lemma}[integration by parts with piecewise data]\label{U:lem:ibp}
If $p,q\in C^1([c,d])$ and $p',q'$ are Lipschitz, then $\int_c^dp\,q''=[p\,q'-p'q]_c^d+\int_c^dp''q$ (with $p''$, $q''$ the a.e.\ derivatives, bounded).
\end{lemma}
\begin{proof}\tH\ $p\,q'-p'q$ is Lipschitz with derivative $p\,q''-p''q$ a.e. (\tS\ S7 for the algebra); integrate.
\end{proof}

\subsection{Proof of Theorem~\ref{U:thm}}
\emph{(a)} $P\in C^1$ and $\mathcal A\in C^\infty$ by Proposition~\ref{U:prop:mod}, so $(a,b,\beta)=\mathcal A(P(\tau))$ is $C^1$. $u=\mathcal U(P(\tau);x)$ with $\mathcal U\in C^\infty$ (Lemma~\ref{U:lem:graph}): the chain rule gives continuity of $\partial_x^ku$ and of $\partial_\tau\partial_x^ku=\nabla_P\partial_x^k\mathcal U\cdot P'$, which equals $\partial_x^k\partial_\tau u$. For $w=u-\alpha\hat\psi_3-\alpha^2\hat h_2$: $\hat\psi_3,\hat h_2\in C^4$, piecewise $C^\infty$ with bounded fifth derivative off $B_0$, and $\partial_\tau(\alpha\hat\psi_3^{(k)})=\alpha'\hat\psi_3^{(k)}$. Products with $\zeta,\chi\in C^3$ (bounded $\zeta'''',\chi''''$ off their break points) give the stated classes. The $w$-equation is the graph equation \tR\ minus the profile part, an identity between functions that are continuous off $B$ (it contains $u''''$, $\hat\psi_3''''$, continuous); differentiating it in $x$ off $B$ and multiplying by $\zeta$ gives (T1)--(T12), whose derivation in \fn{EI_SPEC.md} \S2.3 is pointwise algebra \tH\tS.

\emph{(b)} Apply Lemma~\ref{U:lem:leib} on the compact supports ($|x|\le34.25$ for $\tilde w$; $[10.25,60]$ for $v$) with $\phi=\tilde w^{(j)},v^{(j)}$ ($j\le2$): $\partial_\tau\tilde w^{(j)}=\sum\binom ji\chi^{(j-i)}\partial_\tau w^{(i)}$ is continuous by (a). Then Lemma~\ref{U:lem:ibp} with $p=Wv''$ (so $p'=(Wv'')'$ is Lipschitz since $v'''$ is and $W\in C^4$) and $q=v_\tau=\zeta g_\tau$ ($g_\tau=\partial_xw_\tau\in C^3$ in $x$ since $w_\tau=u_\tau-\alpha'\hat\psi_3-(\alpha^2)'\hat h_2$ is $C^4$ in $x$), on $[10.25,60]$: the boundary terms contain $v_\tau,\partial_xv_\tau$ and $v'',v'''$, which vanish at $10.25$ and $60$ because $\zeta,\zeta',\zeta'',\zeta'''$ do. So $\int Wv''\partial_\tau v''=\int(Wv'')''\,v_\tau=\ip{D_2v}{v_\tau}_W$; the formula for $D_2v$ is S7. For $B^2$: $\tilde w$ is supported in $|x|\le34.25$, $\tilde w'''$ is Lipschitz, $\rho$ is smooth, so the same computation on $[-34.25,34.25]$ has zero boundary terms; on the half line, the boundary terms at $x=0$ vanish by evenness ($\tilde w',\tilde w''',\rho'$ odd). The level-0 identity and the modulation-lemma integrations by parts are Lemma~\ref{U:lem:ibp} (first-order version) with $\zeta$, $\chi$ vanishing to third order at the ends. $\Lambda=E+\mu E_I$ is then $C^1$, in particular absolutely continuous.

\emph{(c)} On $\mathrm{OR}$, $\hat\psi_3,\hat h_2$ are the smooth extensions ($x\ge28$), so $g=w'$ is $C^\infty$ in $x$ and $\partial_x^lg$ is jointly continuous for every $l$. $M_l^2=\int_{\mathrm{OR}}(\partial^lg)^2$, $N_l^2=\int_{\mathrm{OR}}W(\partial^lg)^2$ and $s_2=\max_{\mathrm{OR}}|g''|$ are continuous in $\tau$ (uniform continuity on the compact set). $K^{\rm OR}_0,K^{\rm OR}_2$ are integrals over $\mathrm{OR}$ of products of these with $\zeta^{(i)}$ (bounded, $\zeta''''$ with jumps only at $57,60$) and with coefficient functions continuous in $(\tau,x)$ (they depend on $u_0,w$ and their $x$-derivatives and on $a,b$): continuous by dominated convergence. $\varphi_2$ is a polynomial in $N_3,N_4,s_2$ with fixed coefficients (\fn{flux2_lemma_v44.tex}, Lemma ``flux lemma v2''(c)): continuous. Agmon on an interval of length $3$: for $f=g''\in C^1$, choose $x_0$ with $f(x_0)^2\le\frac13\int f^2$; then $f(x)^2=f(x_0)^2+\int_{x_0}^x2ff'\le\frac13M_2^2+2M_2M_3$ (S8). (F:eq:star) is the $E_I$ inequality of the verifier with $K^{\rm OR}$ bounded by the flux lemma; it is an inequality between continuous functions of $\tau$ obtained from the identities of (b), which hold at every $\tau$.

\emph{(d)} By Lemma~\ref{U:lem:graph}, $f(\tau,y)=\partial_y\mathcal H(P(\tau);y)$ is $C^\infty$ in $y$ with $\partial_\tau\partial_y^kf$ continuous, and $U(\alpha(\tau),y)=1.3056\,\alpha(\tau)y^{1.0246}$ is $C^\infty$ in $y>0$, $C^1$ in $\tau$. So $z^A_k=\partial_s^k(f-U)$ ($s=\ln y$) has continuous $\partial_\tau z^A_k$ on $[\tau_1,\tau']\times[\ln14,\ln700]$. Lemma~\ref{U:lem:leib} with $\Theta=\omega\eta_A^2$ gives $E^A_j\in C^1$. Every integration by parts in \fn{lemmaS_A_v33.tex} \S2 is of a smooth $z$-product against a coefficient that is continuous and piecewise $C^1$ (with $\eta_{\rm in},\eta_o\in C^2$, $S'''$ jumping at the ends of the ramps); on each piece Lemma~\ref{U:lem:ibp} applies, and the interior point contributions are exactly the point terms of its Lemma~P \tI. $X_k(\tau_a,T)=\int_{\tau_a}^{\tau_a+T}x_k$ with $x_k(\tau)=\int\varepsilon\omega(z^A_{k+2})^2$ continuous, so $X_k$ is continuous in $(\tau_a,T)$; likewise for $V_m^2$ (continuous by (b) and dominated convergence for $V_4$). The last sentence is Lemma~\ref{U:lem:leib} again (the window weights of Lemma~U depend on $\tau$ through $s-s_\rho(\tau)$, smoothly). \qed

%======================================================================
\section{What is not proved here}\label{U:sec:gaps}
\begin{enumerate}[leftmargin=*]
\item \textbf{(S4.1) smoothness at $t_1$}: \emph{closed in text v38} by Theorem C-1 of \fn{C1_identification_v44.tex} (identification of the step-4 solution with the smooth Lean flow from $\Gamma_0$, uniqueness in $C^0H^5\cap C^1L^2$ re-proved there by Gr\"onwall, P2 in the a-priori form of its Lemma (P2-a)); it inherits the Lean items below and C-1's own gap list (text v42: DKS is no longer an input).
\item \textbf{Lean} (text v42). (L1), (L2), (T1)--(T5), (V), (Sy), (Sh), (Gl), (PC), (MS), (GE) are compiled theorems (\S\ref{U:sec:lean}). Open: (a) a full \fn{lake build CWFlow.Curve.All} (so far each file, the aggregator \fn{All.lean} included, is compiled against the repository's existing build outputs; text v43, delta check fix23 S-2: text v42 said here that \fn{All.lean} was never compiled, contrary to \S\ref{U:sec:lean}); (b) the commit and push of \fn{CWFlow/Curve/} and the citation of that commit, in place of ``commit pending'' with the snapshot SHA256 sums; the two hand bridges of review W7 (S3a, S3b) that text v42 first listed here are now Lean theorems, (PC), (MS), (GE) \LeanSthree, reviewed in \fn{fix23_delta_check/DELTA_CHECK_v42.md} Part~A (faithful).
\item \textbf{DKS}: no longer an input (text v42). The author item A-2 (``maximal solution'') and the re-reading of DKS for (D1) and for the uniform bounds of C-1 are void; (D1), (D2) are replaced by (T2), (T1).
\item \textbf{(GS)}: \emph{closed in text v38}. At $\tau_1$ by Lemma GS-1 of \fn{GS_lemma_v44.tex} \tI\ (certified, \fn{GS_tau1.json}); in the bootstrap as a member of $\mathrm{BS}$ (\fn{bootstrap_v44.tex}), with strict improvement GS-2 (\fn{GS_boot_v4c.json}), openness GS-3 and closedness GS-4. Step~4's hand-off list with GS next to (H1): \fn{step4_handoff_v44.tex} item (4) (text v39; \fn{step4_proof.tex} is not edited, see \fn{STEP4_PATCH_NOTE_v44.md}); \fn{GS_tau1}, \fn{GS_boot_v4c} and \fn{U1_sym} are stages 23, 24, 26 of \fn{FINAL_manifest_v38.json} (29 stages, \fn{be909bbf...}).
\item \textbf{Read, not re-derived:} the graph equation (\fn{graph_eq.py}, \fn{graph_struct.pkl}); the identification of $(m_{ij})$ with $G+E$ and the Neumann bound $\nu\le0.0466$ (\fn{MOD_IBP.md}); the $C^4$ matching of $\hat\psi_3,\hat h_2$ (\fn{c4fun.py}).
\item \textbf{Everything else} (Lemmas~\ref{U:lem:graph}, \ref{U:lem:leib}, \ref{U:lem:ibp}, Proposition~\ref{U:prop:mod}, Theorem~\ref{U:thm}) is \tH, using the \tS\ identities S1--S8. Lemmas~\ref{U:lem:vel}, \ref{U:lem:glue}, \ref{U:lem:max}, \ref{U:lem:sym}, \ref{U:lem:bu} are \tL\ since text v42 (text v43: the second sentence of Lemma~\ref{U:lem:glue}, the piecewise map, is \tH).
\end{enumerate}
\end{document}
